%% ========================================================================
%% BAB 3: Dasar Input dan Output (I/O) pada Linux (Komprehensif)
%% Sumber konsep diperkaya dari OS-Pertemuan 3.
%% Asumsi: environment box dan lstlisting style sudah didefinisikan seperti week02.
%% ========================================================================

\chapter{Dasar Input dan Output (I/O) pada Linux}
\label{ch:basic-io}

\begin{abstract}
Bab ini membahas mekanisme fundamental input/output pada Linux: filosofi Unix ``everything is a file'',
file descriptor, pengalihan I/O standar (\texttt{stdin}, \texttt{stdout}, \texttt{stderr}),
pipa (pipe) dan operator \texttt{|}, pengalihan output ke file (\texttt{>}, \texttt{>>}),
serta \texttt{tee} untuk menampilkan dan menyimpan output sekaligus.
Pembahasan difokuskan untuk kebutuhan praktikum: otomasi, logging, audit sederhana, dan troubleshooting.
\end{abstract}

% ------------------------------------------------------------------------
\section{Filosofi I/O Unix: ``Everything is a File''}
\label{sec:everything-is-file}

\subsection*{Penjelasan Konsep}
Salah satu prinsip kunci Unix/Linux adalah \textbf{``everything is a file''}:
bukan hanya file reguler, tetapi juga perangkat (\texttt{/dev}), informasi proses (\texttt{/proc}),
informasi sistem (\texttt{/sys}), pipe, socket, dan lain-lain diperlakukan sebagai ``file'' yang
bisa diakses melalui antarmuka seragam: \texttt{open()}, \texttt{read()}, \texttt{write()}, \texttt{close()}.

Keuntungan besar dari filosofi ini:
\begin{itemize}
    \item \textbf{Kesederhanaan antarmuka:} cukup memahami satu pola akses I/O.
    \item \textbf{Komposabilitas:} output suatu program bisa jadi input program lain (pipe/redirection).
    \item \textbf{Portabilitas:} detail hardware disembunyikan oleh abstraksi file.
    \item \textbf{Fleksibilitas:} sumber/tujuan I/O bisa diganti tanpa mengubah logika program.
\end{itemize}

\begin{tipbox}
Cara berpikir yang berguna: ``terminal itu file'', ``keyboard itu file'', ``log itu file''.
Makanya kita bisa mengalihkan output ke mana pun: file, pipe, atau bahkan ``lubang hitam'' (\texttt{/dev/null}).
\end{tipbox}

\subsection*{Praktikum 3.0 --- Membaca Informasi Sistem seperti Membaca File (Mudah)}
\textbf{Tujuan:} membuktikan bahwa informasi sistem dapat dibaca seperti file biasa.

\begin{enumerate}
    \item Lihat informasi CPU:
\begin{lstlisting}[language=bash, caption={Membaca info CPU dari /proc}]
cat /proc/cpuinfo | head
\end{lstlisting}

    \item Lihat informasi memori:
\begin{lstlisting}[language=bash, caption={Membaca info memori dari /proc}]
cat /proc/meminfo | head
\end{lstlisting}
\end{enumerate}

\begin{examplebox}[Refleksi]
Walaupun \texttt{/proc/cpuinfo} terlihat seperti file teks, sebenarnya itu adalah \textit{virtual file} yang disediakan kernel.
Ini contoh nyata ``everything is a file''.
\end{examplebox}

% ------------------------------------------------------------------------
\section{File Descriptor di Linux}
\label{sec:fd}

\subsection*{Penjelasan Konsep}
\textbf{File descriptor (FD)} adalah bilangan bulat non-negatif yang menjadi \textit{handle} bagi proses untuk mengakses file/stream:
file biasa, terminal, pipe, socket, dan sebagainya.
FD adalah implementasi konkret dari filosofi ``everything is a file''.

Mengapa FD diperlukan?
\begin{itemize}
    \item \textbf{Abstraksi:} program tidak perlu tahu detail device/filesystem.
    \item \textbf{Keamanan:} kernel mengontrol akses via FD (izin baca/tulis).
    \item \textbf{Efisiensi:} integer murah untuk dipass dan dikelola.
    \item \textbf{Multiplexing:} satu proses dapat membuka banyak file sekaligus (FD 3,4,5,...).
\end{itemize}

\begin{notebox}
Secara internal, kernel menyimpan tabel FD per proses.
Saat \texttt{open()} dipanggil, kernel memberi FD terkecil yang tersedia (setelah 0,1,2 biasanya mulai dari 3).
\end{notebox}

\subsection*{File Descriptor Standar}
Setiap proses yang berjalan dari terminal umumnya memiliki tiga FD standar:

\begin{center}
\begin{tabular}{|c|c|l|}
\hline
FD & Nama & Fungsi \\
\hline
0 & stdin  & Input standar (keyboard/terminal) \\
1 & stdout & Output normal program \\
2 & stderr & Pesan error/diagnostik \\
\hline
\end{tabular}
\end{center}

\subsection*{Praktikum 3.1 --- Melihat FD Shell Aktif (Mudah)}
\textbf{Tujuan:} melihat FD 0,1,2 pada shell saat ini.

\begin{enumerate}
    \item Tampilkan PID shell:
\begin{lstlisting}[language=bash, caption={PID shell}]
echo $$
\end{lstlisting}

    \item Lihat FD yang terbuka:
\begin{lstlisting}[language=bash, caption={FD shell pada /proc}]
ls -l /proc/$$/fd
\end{lstlisting}

    \item Amati FD 0,1,2 mengarah ke \texttt{/dev/pts/X} (pseudo-terminal).
\end{enumerate}

\begin{exercisebox}[3.1]
Tuliskan hasil mapping FD 0,1,2 pada sistem Anda, lalu jelaskan mengapa semuanya mengarah ke \texttt{/dev/pts/X}.
\end{exercisebox}

\subsection*{Praktikum 3.2 --- FD untuk Logging Terpisah (Riil)}
\textbf{Skenario:} memisahkan output normal dan error untuk debugging cepat.

\begin{enumerate}
    \item Siapkan folder:
\begin{lstlisting}[language=bash, caption={Folder log praktikum}]
mkdir -p ~/log-praktikum
\end{lstlisting}

    \item Jalankan perintah campuran dan pisahkan:
\begin{lstlisting}[language=bash, caption={Pisahkan stdout dan stderr}]
(ls /etc; ls /folder_tidak_ada) > ~/log-praktikum/stdout.log 2> ~/log-praktikum/stderr.log
\end{lstlisting}

    \item Baca log:
\begin{lstlisting}[language=bash, caption={Baca hasil log}]
cat ~/log-praktikum/stdout.log
cat ~/log-praktikum/stderr.log
\end{lstlisting}
\end{enumerate}

% ------------------------------------------------------------------------
\section{Pengalihan I/O Standar (stdin, stdout, stderr)}
\label{sec:redirection}

\subsection*{Penjelasan Konsep}
\textbf{Redirection} adalah mekanisme shell untuk mengubah sumber input atau tujuan output.
Ini penting untuk otomasi, logging, batch processing, dan membangun workflow data.

\subsubsection*{A. Pengalihan stdin (\texttt{<})}
Secara default, program membaca input dari keyboard. Dengan \texttt{<}, input bisa berasal dari file.

\begin{lstlisting}[language=bash, caption={Contoh redirect stdin dari file}]
sort < file.txt
wc -l < /etc/passwd
\end{lstlisting}

\begin{tipbox}
Redirect stdin berguna untuk input besar, input berulang (testing), dan otomasi tanpa interaksi user.
\end{tipbox}

\subsubsection*{B. Pengalihan stdout (\texttt{>} dan \texttt{>>})}
Secara default stdout tampil ke terminal. Dengan \texttt{>} output disimpan ke file (overwrite),
dengan \texttt{>>} output ditambahkan di akhir file (append).

\begin{warningbox}
\texttt{>} menimpa file tanpa peringatan.
Untuk file log, hampir selalu lebih aman memakai \texttt{>>}.
\end{warningbox}

\subsubsection*{C. Pengalihan stderr (\texttt{2>} dan \texttt{2>>})}
stderr adalah channel terpisah untuk error/diagnostik agar data (stdout) tidak tercampur error.
Dengan \texttt{2>} error disimpan ke file, dengan \texttt{2>>} error di-append.

\begin{lstlisting}[language=bash, caption={Redirect stderr}]
find / -name "*.conf" 2> errors.txt
grep "pattern" file.txt 2> /dev/null
\end{lstlisting}

\subsubsection*{D. Menggabungkan stdout dan stderr (Perhatikan Urutan!)}
Ada beberapa metode umum:

\begin{lstlisting}[language=bash, caption={Metode umum redirect stdout+stderr}]
# File berbeda
command > output.txt 2> error.txt

# Satu file (metode klasik) - URUTAN penting
command > all-output.txt 2>&1

# Syntax modern bash (Bash 4+)
command &> all-output.txt

# Append keduanya ke satu file
command >> all-output.txt 2>&1
\end{lstlisting}

\begin{warningbox}
Kesalahan umum: \texttt{command 2>\&1 > file} (urutan terbalik).
Ini membuat stderr tetap ke terminal karena saat \texttt{2>\&1} dieksekusi, stdout masih ke terminal.
Urutan benar: \texttt{> file 2>\&1}.
\end{warningbox}

\subsubsection*{E. Here Document (Input Multi-baris tanpa File Sementara)}
Here document (heredoc) memungkinkan memberi input multi-baris langsung pada command/script.

\begin{lstlisting}[language=bash, caption={Here document untuk input multi-baris}]
cat << EOF
Baris 1
Baris 2
EOF

# Mencegah ekspansi variabel dengan delimiter bertanda kutip
cat << 'EOF' > config.txt
server=localhost
port=8080
EOF
\end{lstlisting}

\subsection*{Praktikum 3.3 --- Redirection Dasar (Mudah)}
\textbf{Tujuan:} menguasai \texttt{>}, \texttt{>>}, \texttt{2>}.

\begin{enumerate}
    \item Simpan daftar file:
\begin{lstlisting}[language=bash, caption={Redirect stdout}]
ls > daftar_file.txt
\end{lstlisting}

    \item Append timestamp:
\begin{lstlisting}[language=bash, caption={Append output}]
date >> daftar_file.txt
\end{lstlisting}

    \item Redirect error:
\begin{lstlisting}[language=bash, caption={Redirect stderr}]
ls file_tidak_ada 2> error.txt
\end{lstlisting}

    \item Verifikasi:
\begin{lstlisting}[language=bash, caption={Cek hasil redirection}]
cat daftar_file.txt
cat error.txt
\end{lstlisting}
\end{enumerate}

\subsection*{Praktikum 3.4 --- Noclobber untuk Mencegah Overwrite (Riil Safety)}
\textbf{Skenario:} mencegah file tertimpa tak sengaja (praktik aman di server).

\begin{enumerate}
    \item Aktifkan noclobber:
\begin{lstlisting}[language=bash, caption={Aktifkan noclobber}]
set -o noclobber
\end{lstlisting}

    \item Coba overwrite file yang sudah ada (harus gagal jika file ada):
\begin{lstlisting}[language=bash, caption={Overwrite diblokir noclobber}]
echo "data" > daftar_file.txt
\end{lstlisting}

    \item Paksa overwrite dengan \texttt{>|}:
\begin{lstlisting}[language=bash, caption={Paksa overwrite}]
echo "data paksa" >| daftar_file.txt
\end{lstlisting}

    \item Nonaktifkan noclobber:
\begin{lstlisting}[language=bash, caption={Nonaktifkan noclobber}]
set +o noclobber
\end{lstlisting}
\end{enumerate}

\subsection*{Praktikum 3.5 --- /dev/null untuk ``Membuang'' Output (Riil)}
\textbf{Skenario:} membuang error yang diharapkan (mis. permission denied saat scanning sistem).

\begin{enumerate}
    \item Jalankan pencarian dan buang error:
\begin{lstlisting}[language=bash, caption={Buang stderr ke /dev/null}]
find / -name "*.conf" 2> /dev/null | head
\end{lstlisting}

    \item Buang semua output (stdout+stderr):
\begin{lstlisting}[language=bash, caption={Buang semua output}]
command > /dev/null 2>&1
\end{lstlisting}
\end{enumerate}

\begin{warningbox}
Hati-hati: sekali output masuk \texttt{/dev/null}, hilang permanen.
Pastikan Anda memang tidak memerlukan output tersebut.
\end{warningbox}

% ------------------------------------------------------------------------
\section{Pipa dan Operator \texttt{|}}
\label{sec:pipes}

\subsection*{Penjelasan Konsep}
\textbf{Pipe} adalah mekanisme IPC (inter-process communication) yang menghubungkan stdout suatu proses menjadi stdin proses lain
dengan operator \texttt{|}. Pipe mengimplementasikan Unix Philosophy:
\textit{buat program kecil yang melakukan satu hal dengan baik, lalu gabungkan.}

Manfaat pipe:
\begin{itemize}
    \item \textbf{Efisiensi:} data mengalir antar proses tanpa file perantara (minim I/O disk).
    \item \textbf{Paralelisme:} tahapan pipeline dapat berjalan \textit{concurrent}.
    \item \textbf{Streaming:} cocok untuk data besar karena diproses bertahap (chunk-by-chunk).
\end{itemize}

\begin{notebox}
Pipe memiliki buffer di kernel (umumnya puluhan KB).
Jika buffer penuh, proses penulis akan \textit{blocked} sampai proses pembaca mengonsumsi data.
Ini menjadi flow control alami.
\end{notebox}

\subsection*{Praktikum 3.6 --- Pipeline Dasar (Mudah)}
\begin{enumerate}
    \item Hitung jumlah item:
\begin{lstlisting}[language=bash, caption={Hitung item di direktori}]
ls | wc -l
\end{lstlisting}

    \item Filter file \texttt{.txt}:
\begin{lstlisting}[language=bash, caption={Filter txt}]
ls | grep ".txt"
\end{lstlisting}

    \item Pipeline bertingkat:
\begin{lstlisting}[language=bash, caption={Filter lalu hitung}]
ls -l | grep ".txt" | wc -l
\end{lstlisting}
\end{enumerate}

\subsection*{Praktikum 3.7 --- Pipeline ``Kondisi Riil'' (Analisis Log/Proses)}
\textbf{Skenario A (Proses):} cari proses tertentu lalu simpan hasilnya.
\begin{lstlisting}[language=bash, caption={Filter proses dan simpan}]
ps aux | grep ssh | sort -k3 -rn | head -10 > top-ssh-processes.txt
\end{lstlisting}

\textbf{Skenario B (Log autentikasi):} investigasi login gagal.
\begin{lstlisting}[language=bash, caption={Analisis login gagal}]
sudo cat /var/log/auth.log | grep "Failed" | tail -n 10
sudo cat /var/log/auth.log | grep "Failed" | wc -l
\end{lstlisting}

\subsection*{Tambahan Konsep: Named Pipe (FIFO)}
Selain pipe ``sementara'' dari operator \texttt{|}, Linux juga mendukung \textbf{named pipe} (FIFO) yang punya nama di filesystem
dan bisa dipakai antar terminal/proses terpisah.

\subsection*{Praktikum 3.8 --- Named Pipe (Riil IPC Sederhana)}
\textbf{Tujuan:} simulasi komunikasi antar proses/terminal.

\begin{enumerate}
    \item Buat FIFO:
\begin{lstlisting}[language=bash, caption={Membuat FIFO}]
mkfifo mypipe
\end{lstlisting}

    \item \textbf{Terminal 1} (menulis):
\begin{lstlisting}[language=bash, caption={Menulis ke FIFO}]
echo "Hello from Terminal 1" > mypipe
\end{lstlisting}

    \item \textbf{Terminal 2} (membaca):
\begin{lstlisting}[language=bash, caption={Membaca dari FIFO}]
cat < mypipe
\end{lstlisting}

    \item Hapus FIFO:
\begin{lstlisting}[language=bash, caption={Menghapus FIFO}]
rm mypipe
\end{lstlisting}
\end{enumerate}

% ------------------------------------------------------------------------
\section{Pengalihan Output ke File (\texttt{>} dan \texttt{>>})}
\label{sec:redirect-to-file}

\subsection*{Penjelasan Konsep}
Pengalihan output ke file adalah dasar pembuatan log dan laporan.

\textbf{Overwrite (\texttt{>})} cocok untuk hasil yang harus selalu fresh.
\textbf{Append (\texttt{>>})} cocok untuk log dan histori.

\begin{tipbox}
Rule of thumb:
\texttt{>} = laporan baru,
\texttt{>>} = log yang terus bertambah.
\end{tipbox}

\subsection*{Praktikum 3.9 --- Overwrite vs Append (Mudah)}
\begin{enumerate}
    \item Overwrite:
\begin{lstlisting}[language=bash, caption={Overwrite}]
echo "Baris Pertama" > latihan.txt
\end{lstlisting}

    \item Append:
\begin{lstlisting}[language=bash, caption={Append}]
echo "Baris Kedua" >> latihan.txt
\end{lstlisting}

    \item Cek:
\begin{lstlisting}[language=bash, caption={Cek isi}]
cat latihan.txt
\end{lstlisting}

    \item Timpa:
\begin{lstlisting}[language=bash, caption={Timpa}]
echo "Ditimpa" > latihan.txt
cat latihan.txt
\end{lstlisting}
\end{enumerate}

\subsection*{Praktikum 3.10 --- Laporan Harian Server (Riil)}
\textbf{Skenario:} membuat daily report manual (dasar automasi cron).

\begin{enumerate}
    \item Header + timestamp:
\begin{lstlisting}[language=bash, caption={Membuat header laporan}]
echo "=== DAILY REPORT ===" > daily_report.txt
date >> daily_report.txt
\end{lstlisting}

    \item Disk + memory:
\begin{lstlisting}[language=bash, caption={Tambahkan disk dan memory}]
echo "--- DISK USAGE ---" >> daily_report.txt
df -h >> daily_report.txt
echo "--- MEMORY ---" >> daily_report.txt
free -h >> daily_report.txt
\end{lstlisting}

    \item User login:
\begin{lstlisting}[language=bash, caption={Tambahkan user login}]
echo "--- LOGGED IN USERS ---" >> daily_report.txt
who >> daily_report.txt
\end{lstlisting}

    \item Tampilkan:
\begin{lstlisting}[language=bash, caption={Tampilkan laporan}]
cat daily_report.txt
\end{lstlisting}
\end{enumerate}

% ------------------------------------------------------------------------
\section{Penggunaan \texttt{tee} untuk Output Ganda}
\label{sec:tee}

\subsection*{Penjelasan Konsep}
Perintah \texttt{tee} membaca dari stdin lalu menulis ke stdout dan ke satu/lebih file.
Nama ``tee'' mengilustrasikan bentuk huruf T: aliran data dibagi dua.

Kapan \texttt{tee} berguna?
\begin{itemize}
    \item \textbf{Real-time monitoring + logging:} lihat proses sambil simpan log.
    \item \textbf{Audit trail:} dokumentasi langkah maintenance/deployment.
    \item \textbf{Debug pipeline:} ``menyadap'' data intermediate tanpa mengubah alur.
    \item \textbf{Menulis ke file root:} gunakan \texttt{sudo tee} (redirection biasa tidak ikut sudo).
\end{itemize}

\begin{notebox}
Redirection \texttt{>} diproses oleh shell \textit{sebelum} \texttt{sudo}.
Karena itu \texttt{sudo echo "x" > /etc/file} sering gagal.
Solusinya: \texttt{echo "x" | sudo tee /etc/file}.
\end{notebox}

\subsection*{Praktikum 3.11 --- tee Dasar (Mudah)}
\begin{enumerate}
    \item Tampilkan dan simpan:
\begin{lstlisting}[language=bash, caption={tee dasar}]
ls | tee hasil.txt
\end{lstlisting}

    \item Append:
\begin{lstlisting}[language=bash, caption={tee append}]
date | tee -a hasil.txt
\end{lstlisting}

    \item Cek:
\begin{lstlisting}[language=bash, caption={Cek hasil}]
cat hasil.txt
\end{lstlisting}
\end{enumerate}

\subsection*{Praktikum 3.12 --- tee untuk Logging APT Update (Riil)}
\textbf{Skenario:} saat update server, Anda ingin melihat progress dan punya log.

\begin{enumerate}
    \item Jalankan update dan simpan log:
\begin{lstlisting}[language=bash, caption={apt update dengan tee}]
sudo apt update | tee apt-update.log
\end{lstlisting}

    \item (Opsional) upgrade dan simpan log:
\begin{lstlisting}[language=bash, caption={apt upgrade dengan tee}]
sudo apt upgrade -y | tee apt-upgrade.log
\end{lstlisting}
\end{enumerate}

\subsection*{Praktikum 3.13 --- tee untuk Menulis File Butuh Root (Riil)}
\textbf{Skenario:} menambah entri ke \texttt{/etc/hosts}.

\begin{enumerate}
    \item Append aman dengan \texttt{sudo tee -a}:
\begin{lstlisting}[language=bash, caption={Append ke /etc/hosts}]
echo "127.0.0.1 contoh.local" | sudo tee -a /etc/hosts
\end{lstlisting}

    \item Verifikasi:
\begin{lstlisting}[language=bash, caption={Cek hosts}]
tail -n 3 /etc/hosts
\end{lstlisting}
\end{enumerate}

\subsection*{Catatan Penting: Exit Code pada Pipeline (\texttt{\$?} vs \texttt{PIPESTATUS})}
Pada pipeline, \texttt{\$?} hanya menyimpan exit code command \textit{terakhir}.
Untuk mengambil exit code tiap tahapan pipeline, gunakan array \texttt{PIPESTATUS}.

\begin{lstlisting}[language=bash, caption={PIPESTATUS untuk cek exit code pipeline}]
command | tee output.txt
echo ${PIPESTATUS[0]}   # exit code command
echo ${PIPESTATUS[1]}   # exit code tee
\end{lstlisting}

\begin{warningbox}
Jika Anda membuat script logging menggunakan \texttt{tee}, jangan hanya mengandalkan \texttt{\$?}.
Gunakan \texttt{PIPESTATUS} agar tahu apakah command utama sukses/gagal.
\end{warningbox}

% ------------------------------------------------------------------------
\section{Tugas dan Latihan Mandiri}
\label{sec:tasks-io}

\subsection*{A. Latihan Konseptual}
\begin{enumerate}
    \item Jelaskan mengapa Unix memisahkan \texttt{stdout} dan \texttt{stderr}. Berikan satu contoh manfaatnya pada pipeline.
    \item Jelaskan perbedaan \texttt{>}, \texttt{>>}, \texttt{2>}, \texttt{2>>}, \texttt{2>\&1}, dan \texttt{\&>}.
    \item Mengapa urutan \texttt{> file 2>\&1} penting? Jelaskan dengan logika FD.
\end{enumerate}

\subsection*{B. Latihan Praktikum (Mudah)}
\begin{enumerate}
    \item Buat pipeline yang membaca \texttt{/etc/passwd}, mengambil username (kolom pertama), urutkan alfabetis, dan simpan ke \texttt{sorted-users.txt}.
    \item Buat perintah yang mencari semua file \texttt{.conf} di sistem tetapi membuang pesan ``Permission denied'', lalu hitung jumlah file yang ditemukan.
\end{enumerate}

\subsection*{C. Latihan Praktikum (Kondisi Riil)}
\begin{enumerate}
    \item Buat script/perintah yang menampilkan 10 file terbesar di \texttt{/var/log}, tampilkan di terminal \textbf{dan} simpan ke \texttt{large-logs.txt} menggunakan \texttt{tee}. Error disimpan ke \texttt{error.log}.
    \item Buat ``monitoring mini'' selama 1 menit (12 iterasi) yang setiap 5 detik mencatat ringkasan CPU dan memory (boleh minimal: \texttt{uptime} dan \texttt{free -h}),
    tampilkan dan simpan ke file log dengan timestamp.
\end{enumerate}

\begin{tipbox}
Jika sudah nyaman dengan bab ini, Anda siap naik tingkat ke: redirection lanjutan, scripting bash yang robust, dan automasi dengan \texttt{cron} atau \texttt{systemd timer}.
\end{tipbox}